\documentclass{amsart}
% For dealing with references we use the comment environment
\usepackage{verbatim}
\newenvironment{reference}{\comment}{\endcomment}
%\newenvironment{reference}{}{}
\newenvironment{slogan}{\comment}{\endcomment}
\newenvironment{history}{\comment}{\endcomment}

% For commutative diagrams we use Xy-pic
\usepackage[all]{xy}

% We use 2cell for 2-commutative diagrams.
\xyoption{2cell}
\UseAllTwocells

% We use multicol for the list of chapters between chapters
\usepackage{multicol}

% This is generally recommended for better output
\usepackage{lmodern}
\usepackage[T1]{fontenc}

% For cross-file-references

% Package for hypertext links:
\usepackage{hyperref}

% For any local file, say "hello.tex" you want to link to please

% Theorem environments.
%
\theoremstyle{plain}
\newtheorem{theorem}[subsection]{Theorem}
\newtheorem{proposition}[subsection]{Proposition}
\newtheorem{lemma}[subsection]{Lemma}

\theoremstyle{definition}
\newtheorem{definition}[subsection]{Definition}
\newtheorem{example}[subsection]{Example}
\newtheorem{exercise}[subsection]{Exercise}
\newtheorem{situation}[subsection]{Situation}

\theoremstyle{remark}
\newtheorem{remark}[subsection]{Remark}
\newtheorem{remarks}[subsection]{Remarks}

\numberwithin{equation}{subsection}

% Macros
%
\def\lim{\mathop{\mathrm{lim}}\nolimits}
\def\colim{\mathop{\mathrm{colim}}\nolimits}
\def\Spec{\mathop{\mathrm{Spec}}}
\def\Hom{\mathop{\mathrm{Hom}}\nolimits}
\def\Ext{\mathop{\mathrm{Ext}}\nolimits}
\def\SheafHom{\mathop{\mathcal{H}\!\mathit{om}}\nolimits}
\def\SheafExt{\mathop{\mathcal{E}\!\mathit{xt}}\nolimits}
\def\Sch{\mathit{Sch}}
\def\Mor{\mathop{\mathrm{Mor}}\nolimits}
\def\Ob{\mathop{\mathrm{Ob}}\nolimits}
\def\Sh{\mathop{\mathit{Sh}}\nolimits}
\def\NL{\mathop{N\!L}\nolimits}
\def\CH{\mathop{\mathrm{CH}}\nolimits}
\def\proetale{{pro\text{-}\acute{e}tale}}
\def\etale{{\acute{e}tale}}
\def\QCoh{\mathit{QCoh}}
\def\Ker{\mathop{\mathrm{Ker}}}
\def\Im{\mathop{\mathrm{Im}}}
\def\Coker{\mathop{\mathrm{Coker}}}
\def\Coim{\mathop{\mathrm{Coim}}}

% Boxtimes
%
\DeclareMathSymbol{\boxtimes}{\mathbin}{AMSa}{"02}

%
% Macros for moduli stacks/spaces
%
\def\QCohstack{\mathcal{QC}\!\mathit{oh}}
\def\Cohstack{\mathcal{C}\!\mathit{oh}}
\def\Spacesstack{\mathcal{S}\!\mathit{paces}}
\def\Quotfunctor{\mathrm{Quot}}
\def\Hilbfunctor{\mathrm{Hilb}}
\def\Curvesstack{\mathcal{C}\!\mathit{urves}}
\def\Polarizedstack{\mathcal{P}\!\mathit{olarized}}
\def\Complexesstack{\mathcal{C}\!\mathit{omplexes}}
% \Pic is the operator that assigns to X its picard group, usage \Pic(X)
% \Picardstack_{X/B} denotes the Picard stack of X over B
% \Picardfunctor_{X/B} denotes the Picard functor of X over B
\def\Pic{\mathop{\mathrm{Pic}}\nolimits}
\def\Picardstack{\mathcal{P}\!\mathit{ic}}
\def\Picardfunctor{\mathrm{Pic}}
\def\Deformationcategory{\mathcal{D}\!\mathit{ef}}

\setlength{\emergencystretch}{3em}
\begin{document}
\title[Stacks Project, Tag 0G2H]{264 --- Deligne systems represent derived Hom after restriction to an open subscheme}
\maketitle
\noindent Copyright (C) 2005--2025 Johan de Jong. Stacks Project authors. Source: \href{https://stacks.math.columbia.edu/tag/0G2H}{Tag 0G2H}.

\noindent Editorial difficulty note (not author text). Difficulty H2: The proof reduces to Ext of coherent sheaves, then proves both injectivity and surjectivity by induction on Ext degree. Artin-Rees controls the changing kernels and powers of the boundary ideal while coherent extensions kill and lift extension classes. The pro-system comparison is a long exactness-and-lifting argument..

\noindent Editorial provenance note (not author text). History: excerpt from revision \texttt{a0444\allowbreak 6e57e\allowbreak c1fbc\allowbreak 252a8\allowbreak 71afc\allowbreak ec775\allowbreak 2fb28\allowbreak 07b14}, duality.tex, lines 7702--7869. Reference presentation was adapted; original-author context and core support blocks were retained or added. The mathematical wording is unchanged. Licensed under GNU FDL 1.2 or later; no invariant sections or cover texts. See ../licenses/COPYING. Context blocks reproduce author definitions and supporting results; they are not additional counted problems.

\section{Extension by zero for coherent modules}
\label{section-extension-by-zero}

\noindent
The material in this section and the next few can be found
in the appendix by Deligne of \href{https://stacks.math.columbia.edu/bibliography}{Bibliography: RD}.

\medskip\noindent
In this section $j : U \to X$ will be an open immersion of Noetherian schemes.
We are going to consider inverse systems $(K_n)$ in
$D^b_{\textit{Coh}}(\mathcal{O}_X)$ constructed as follows.
Let $\mathcal{F}^\bullet$ be a bounded complex of coherent
$\mathcal{O}_X$-modules. Let $\mathcal{I} \subset \mathcal{O}_X$
be a quasi-coherent sheaf of ideals with $V(\mathcal{I}) = X \setminus U$.
Then we can set
$$
K_n = \mathcal{I}^n\mathcal{F}^\bullet
$$
More precisely, $K_n$ is the object of $D^b_{\textit{Coh}}(\mathcal{O}_X)$
represented by the complex whose term in degree $q$
is the coherent submodule $\mathcal{I}^n\mathcal{F}^q$ of $\mathcal{F}^q$.
Observe that the maps $\ldots \to K_3 \to K_2 \to K_1$ induce isomorphisms
on restriction to $U$. Let us call such a system a {\it Deligne system}.



\begin{lemma}
\label{lemma-lift-map}
Let $j : U \to X$ be an open immersion of Noetherian schemes.
Let $(K_n)$ be a Deligne system and denote
$K \in D^b_{\textit{Coh}}(\mathcal{O}_U)$ the value
of the constant system $(K_n|_U)$. Let $L$ be an object of
$D^b_{\textit{Coh}}(\mathcal{O}_X)$.
Then $\colim \Hom_X(K_n, L) = \Hom_U(K, L|_U)$.
\end{lemma}

\begin{proof}
Let $L \to M \to N \to L[1]$ be a distinguished triangle in
$D^b_{\textit{Coh}}(\mathcal{O}_X)$. Then we obtain
a commutative diagram
$$
\xymatrix{
\ldots \ar[r] &
\colim \Hom_X(K_n, L) \ar[r] \ar[d] &
\colim \Hom_X(K_n, M) \ar[r] \ar[d] &
\colim \Hom_X(K_n, N) \ar[r] \ar[d] &
\ldots \\
\ldots \ar[r] &
\Hom_U(K, L|_U) \ar[r] &
\Hom_U(K, M|_U) \ar[r] &
\Hom_U(K, N|_U) \ar[r] &
\ldots
}
$$
whose rows are exact by Derived Categories, Lemma
\href{https://stacks.math.columbia.edu/tag/0149}{lemma representable homological (Tag 0149)} and
Algebra, Lemma \href{https://stacks.math.columbia.edu/tag/00DB}{lemma directed colimit exact (Tag 00DB)}.
Hence if the statement of the lemma holds for
$N[-1]$, $L$, $N$, and $L[1]$ then it holds for $M$ by the 5-lemma.
Thus, using the distinguished triangles
for the canonical truncations of $L$ (see Derived Categories, Remark
\href{https://stacks.math.columbia.edu/tag/08J5}{remark truncation distinguished triangle (Tag 08J5)})
we reduce to the case that $L$ has only one nonzero cohomology sheaf.

\medskip\noindent
Choose a bounded complex $\mathcal{F}^\bullet$ of coherent
$\mathcal{O}_X$-modules and a quasi-coherent ideal
$\mathcal{I} \subset \mathcal{O}_X$ cutting out $X \setminus U$
such that $K_n$ is represented by $\mathcal{I}^n\mathcal{F}^\bullet$.
Using ``stupid'' truncations we obtain compatible termwise split short
exact sequences of complexes
$$
0 \to \sigma_{\geq a + 1} \mathcal{I}^n\mathcal{F}^\bullet \to
\mathcal{I}^n\mathcal{F}^\bullet \to
\sigma_{\leq a} \mathcal{I}^n\mathcal{F}^\bullet \to 0
$$
which in turn correspond to compatible systems of distinguished
triangles in $D^b_{\textit{Coh}}(\mathcal{O}_X)$.
Arguing as above we reduce to the case where $\mathcal{F}^\bullet$
has only one nonzero term. This reduces us to the case discussed in
the next paragraph.

\medskip\noindent
Given a coherent $\mathcal{O}_X$-module $\mathcal{F}$ and a coherent
$\mathcal{O}_X$-module $\mathcal{G}$ we
have to show that the canonical map
$$
\colim \Ext^i_X(\mathcal{I}^n\mathcal{F}, \mathcal{G})
\longrightarrow
\Ext^i_U(\mathcal{F}|_U, \mathcal{G}|_U)
$$
is an isomorphism for all $i \geq 0$. For $i = 0$ this is
Cohomology of Schemes, Lemma \href{https://stacks.math.columbia.edu/tag/01YB}{lemma homs over open (Tag 01YB)}.
Assume $i > 0$.

\medskip\noindent
Injectivity. Let $\xi \in \Ext^i_X(\mathcal{I}^n\mathcal{F}, \mathcal{G})$
be an element whose restriction to $U$ is zero. We have to show there exists
an $m \geq n$ such that the restriction of $\xi$ to
$\mathcal{I}^m\mathcal{F} = \mathcal{I}^{m - n}\mathcal{I}^n\mathcal{F}$
is zero. After replacing $\mathcal{F}$ by $\mathcal{I}^n\mathcal{F}$
we may assume $n = 0$, i.e., we have
$\xi \in \Ext^i_X(\mathcal{F}, \mathcal{G})$
whose restriction to $U$ is zero. By
Derived Categories of Schemes, Proposition \href{https://stacks.math.columbia.edu/tag/0FDB}{proposition DCoh (Tag 0FDB)}
we have
$D^b_{\textit{Coh}}(\mathcal{O}_X) = D^b(\textit{Coh}(\mathcal{O}_X))$.
Hence we can compute the $\Ext$ group in the abelian category of
coherent $\mathcal{O}_X$-modules. This implies there exists an
surjection $\alpha : \mathcal{F}'' \to \mathcal{F}$ such that
$\xi \circ \alpha = 0$ (this is where we use that $i > 0$).
Set $\mathcal{F}' = \Ker(\alpha)$ so that we have a short exact
sequence
$$
0 \to \mathcal{F}' \to \mathcal{F}'' \to \mathcal{F} \to 0
$$
It follows that $\xi$ is the image of an element
$\xi' \in \Ext^{i - 1}_X(\mathcal{F}', \mathcal{G})$
whose restriction to $U$ is in the image of
$\Ext^{i - 1}_U(\mathcal{F}''|_U, \mathcal{G}|_U) \to
\Ext^{i - 1}_U(\mathcal{F}'|_U, \mathcal{G}|_U)$.
By Artin-Rees the inverse systems $(\mathcal{I}^n\mathcal{F}')$ and
$(\mathcal{I}^n \mathcal{F}'' \cap \mathcal{F}')$ are pro-isomorphic, see
Cohomology of Schemes, Lemma \href{https://stacks.math.columbia.edu/tag/01YA}{lemma Artin Rees (Tag 01YA)}.
Since we have the compatible system of short exact sequences
$$
0 \to
\mathcal{F}' \cap \mathcal{I}^n\mathcal{F}'' \to
\mathcal{I}^n\mathcal{F}'' \to
\mathcal{I}^n\mathcal{F} \to 0
$$
we obtain a commutative diagram
$$
\xymatrix{
\colim \Ext^{i - 1}_X(\mathcal{I}^n\mathcal{F}'', \mathcal{G})
\ar[r] \ar[d] &
\colim \Ext^{i - 1}_X(\mathcal{F}' \cap \mathcal{I}^n\mathcal{F}'', \mathcal{G})
\ar[r] \ar[d] &
\colim \Ext^i_X(\mathcal{I}^n\mathcal{F}, \mathcal{G})
\ar[d] \\
\Ext^{i - 1}_U(\mathcal{F}''|_U, \mathcal{G}|_U) \ar[r] &
\Ext^{i - 1}_U(\mathcal{F}'|_U, \mathcal{G}|_U) \ar[r] &
\Ext^{i - 1}_U(\mathcal{F}|_U, \mathcal{G}|_U)
}
$$
with exact rows. By induction on $i$ and the comment on inverse systems above
we find that the left two vertical arrows are isomorphisms.
Now $\xi$ gives an element in the top right group which is the image
of $\xi'$ in the middle top group, which in turn maps to an element of
the bottom middle group coming from some element in the left bottom group.
We conclude that $\xi$ maps to zero in
$\Ext^i_X(\mathcal{I}^n\mathcal{F}, \mathcal{G})$
for some $n$ as desired.

\medskip\noindent
Surjectivity. Let $\xi \in \Ext^i_U(\mathcal{F}|_U, \mathcal{G}|_U)$.
Arguing as above using that $i > 0$ we can find an surjection
$\mathcal{H} \to \mathcal{F}|_U$ of coherent $\mathcal{O}_U$-modules
such that $\xi$ maps to zero in $\Ext^i_U(\mathcal{H}, \mathcal{G}|_U)$.
Then we can find a map $\varphi : \mathcal{F}'' \to \mathcal{F}$
of coherent $\mathcal{O}_X$-modules whose restriction to $U$ is
$\mathcal{H} \to \mathcal{F}|_U$, see
Properties, Lemma \href{https://stacks.math.columbia.edu/tag/01PI}{lemma extend finite presentation (Tag 01PI)}.
Observe that the lemma doesn't guarantee $\varphi$ is surjective
but this won't matter (it is possible to pick a surjective $\varphi$
with a little bit of additional work).
Denote $\mathcal{F}' = \Ker(\varphi)$. The short exact sequence
$$
0 \to \mathcal{F}'|_U \to \mathcal{F}''|_U \to \mathcal{F}|_U \to 0
$$
shows that $\xi$ is the image of $\xi'$ in
$\Ext^{i - 1}_U(\mathcal{F}'|_U, \mathcal{G}|_U)$.
By induction on $i$ we can find an $n$ such that
$\xi'$ is the image of some $\xi'_n$ in
$\Ext^{i - 1}_X(\mathcal{I}^n\mathcal{F}', \mathcal{G})$.
By Artin-Rees we can find an $m \geq n$ such that
$\mathcal{F}' \cap \mathcal{I}^m\mathcal{F}'' \subset
\mathcal{I}^n\mathcal{F}'$. Using the short exact sequence
$$
0 \to \mathcal{F}' \cap \mathcal{I}^m\mathcal{F}'' \to
\mathcal{I}^m\mathcal{F}'' \to \mathcal{I}^m\Im(\varphi) \to 0
$$
the image of $\xi'_n$ in
$\Ext^{i - 1}_X(\mathcal{F}' \cap \mathcal{I}^m\mathcal{F}'', \mathcal{G})$
maps by the boundary map to an element $\xi_m$ of
$\Ext^i_X(\mathcal{I}^m\Im(\varphi), \mathcal{G})$ which maps
to $\xi$. Since $\Im(\varphi)$ and $\mathcal{F}$ agree over $U$
we see that $\mathcal{F}/\mathcal{I}^m\Im(\varphi)$ is supported
on $X \setminus U$. Hence there exists an $l \geq m$
such that $\mathcal{I}^l\mathcal{F} \subset \mathcal{I}^m\Im(\varphi)$, see
Cohomology of Schemes, Lemma \href{https://stacks.math.columbia.edu/tag/01Y9}{lemma power ideal kills sheaf (Tag 01Y9)}.
Taking the image of $\xi_m$ in
$\Ext^i_X(\mathcal{I}^l\mathcal{F}, \mathcal{G})$ we win.
\end{proof}
\end{document}
